\section{Kloosterman Sums}
We now turn to a genuinely different family of exponential sums, introduced by Hendrik Kloosterman in 1926 in his refinement of the Hardy--Littlewood circle method for the representation of integers by positive definite quaternary quadratic forms. Instead of twisting two characters of different types, a Kloosterman sum twists a single additive character $\chi$ with the \emph{inversion} map $\gm\mapsto \gm^{-1}$ on $\bbF_q^*$, producing a sum of the form $\sum_{\gm}\chi(a\gm+b\gm^{-1})$.

\begin{Definition}{Kloosterman Sums}{def:kloosterman}
  Let $\chi\in\Xq$ be a non-trivial additive character of $\bbF_q$ and let $a,b\in\bbF_q$. Then the \emph{Kloosterman Sum} associated to $\chi,a,b$ is defined as $$K(\chi;a,b)=\sum_{\gm\in\bbF_q^*}\chi(a\cdot \gm+b\cdot \gm^{-1})$$
\end{Definition}

When one or both of the parameters $a,b$ vanish, the Kloosterman sum collapses to an additive character sum on $\bbF_q^*$, which we have already computed in \AddCharPropsthm. If both $a=b=0$ then for all $\gm\in\bbF_q^*$, $\chi(a\cdot \gm+b\cdot \gm^{-1})=\chi(0)=1$. Hence $K(\chi;0,0)=q-1$. 

If exactly one of $a,b$ is zero then we have either $$K(\chi;a,b)=\sum_{\gm\in\bbF_q^*}\chi(a\cdot \gm)=-1\text{ or }K(\chi;a,b)=\sum_{\gm\in\bbF_q^*}\chi(b\cdot \gm^{-1})=-1$$ So in both cases $K(\chi;a,b)=-1$. So we have the following observations:
\begin{observation}\label{obs:kloosterman-at-zero}
  Let $\chi\in\Xq$ be a non-trivial additive character of $\bbF_q$ and $a,b\in\bbF_q$ such that $ab=0$. Then
  \begin{enumerate}[label=(\roman*)]
    \item $K(\chi;0,0)=q-1$ if both $a=b=0$;
    \item if exactly one of $a,b$ is $0$ then $K(\chi;a,b)=-1$;
  \end{enumerate}
\end{observation}

Now another property to notice is Kloosterman sum is always real since $$\ov{K(\chi;a,b)}=\sum_{\gm\in\bbF_q^*}\ov{\chi}(a\cdot \gm+b\cdot \gm^{-1})= \sum_{\gm\in\bbF_q^*}\chi(-(a\cdot \gm+b\cdot \gm^{-1}))=\sum_{\gm\in\bbF_q^*}\chi(a\cdot (-\gm)+b\cdot (-\gm)^{-1})=K(\chi;a,b)$$
Therefore we have another observation:
\begin{observation}\label{obs:kloosterman-real}
  For all additive characters $\chi\in\Xq$ and $a,b\in\bbF_p$, the Kloosterman Sum is real i.e. $\ov{K(\chi;a,b)}=K(\ov{\chi};a,b)=K(\chi;-a,-b)=K(\chi;a,b)$. 
\end{observation}
Now we can express Kloosterman sums as product of two additive characters like Jacobi sums is for product of multiplicative characters. For any $\chi\in\Xq$ let $\chi=\chi_{c}$ for some $c\in\bbF_q$ as defined in \AddCharDefthm. Then $$K\lt(\chi_c;a,b\rt)=\sum_{\gm\in\bbF_q^*}\chi_c\lt(a\cdot \gm+b\cdot \gm^{-1}\rt)=\sum_{\gm\in\bbF_q^*}\chi_{a\cdot c}\lt(\gm\rt)\cdot \chi_{-b\cdot c}\lt(\gm\rt)\eqqcolon K\lt(\chi_{ac},\ov{\chi}_{bc}\rt)$$ 
\begin{observation}\label{obs:kloosterman-two-char-def}
  The Kloosterman sum can equivalently be viewed as two-additive-character sum: for any $a,b,c\in\bbF_q$ and $\chi_a,\chi_b,\chi_c\in\Xq$ (in the notation of \AddCharDefthm), $$K\lt(\chi_a,\chi_b\rt)\coloneqq K\lt(\chi_1;a,-b\rt)\qquad\text{and}\qquad K\lt(\chi_c;a,-b\rt)=K\lt(\chi_{ac},\chi_{bc}\rt)$$
\end{observation}

When the ground field has odd characteristic there is a useful alternative expression for $K(\chi;a,b)$ as a sum of a single additive character against the quadratic character $\eta$ evaluated on a quadratic polynomial. This representation is often the starting point for asymptotic bounds on Kloosterman sums in arithmetic contexts.
\begin{lemma}{}{lemma:kloosterman-quadratic}
  Let $\chi\in\Xq$ be a non-trivial additive character of $\bbF_q$ with $q$ odd, let $\eta$ be the quadratic character of $\bbF_q$ (extended by $\eta(0)=0$), and let $a,b\in\bbF_q$ be not both zero. Then $$K(\chi;a,b)=\sum_{d\in\bbF_q}\chi(d)\cdot \eta(d^2-4ab)$$
\end{lemma}
\begin{proof}
  If one of $a,b$ is zero then by \obsref[(ii)]{obs:kloosterman-at-zero} we have $K(\chi;a,b)=-1$. And then at the right hand side we have the sum over $\eta(\alpha^2-4ab)=\eta(\alpha^2)=1$ and therefore the sum is of $\chi(\alpha)$ over all non-zero alpha which by \AddCharPropsthm[(i)] gives $q-1=-1$. Therefore in this case the lemma is satisfied. So now assume $a\cdot b\neq 0$. 

  Let for any $\gm\in\bbF_q^*$, denote $d=a\cdot \gm+b\cdot \gm^{-1}$. Now for any $d\in\bbF_q$ let $Z(d)$ is the number of $\gm\in\bbF_q^*$ such that $d=a\cdot \gm+b\cdot \gm^{-1}$. Now for any $d\in\bbF_q$, $$a\cdot \gm+b\cdot \gm^{-1}=d\iff a\cdot \gm^2-d\cdot \gm+b=0$$Therefore if $P_d(X)=a\cdot X^2-d\dot X+b\in\bbF_q[X]$ be a polynomial then $Z(d)=$ number of zeros of $P_d(X)$. Hence $Z(d)\leq 2$. Let $\Dl_d$ is the discriminant of $P_d$ i.e. $\Dl_d=d^2-4ab$. Hence If $\Dl_d=0$ then $Z(d)=1$. If a non-zero square root of $\Dl_d$ exists in $\bbF_q$ then $\Dl_d$ is a quadratic residue i.e. $\eta(\Dl_d)=1$ then $Z(d)=2$ and if it doesn't exists then $\Dl_d$ is a quadratic non-residue i.e. $\eta(\Dl_d)=-1$ then $Z(d)=0$. Therefore together we can write $Z(d)=\eta(\Dl_d)+1$ with the extended definition of $\eta$ by setting $\eta(0)=0$. Therefore we have  $$K(\chi;a,b)=\sum_{d\in\bbF_q}\chi(d)\cdot Z(d)=\sum_{d\in\bbF_q}\chi(d)\cdot\cdot \eta(\Dl_d)=\sum_{d\in\bbF_q}\chi(d)\cdot\cdot \eta(d^2-4ab)$$Hence we have the lemma.
\end{proof}

One of the central structural result for Kloosterman sums asserts that the sequence of lifted sums $K(\chi^{(s)};a,b)$, viewed as a function of $s$, is controlled by just two complex numbers $\omega_1,\omega_2$ -- exactly the reciprocal roots of the associated $L$-function. This is the Kloosterman analogue of the \DHthm. We will use again the \emph{Dirichlet $L$-functions} in \autoref{sec:l-functions} to prove the following theorem
\begin{Theorem}{}{thm:kloosterman-lift}
  Let $\chi\in\Xq$ be a non-trivial additive character of $\bbF_q$ and let $a,b\in\bbF_q$ with $ab\neq 0$. Suppose $\chi$ is lifted to character $\chi^{(r)}$ of the finite extension field $E$ of $\bbF_q$ with $[E\colon \bbF_q]=r$. Then there exist complex numbers $w_1,w_2$ (depending only on $\chi,a,b$), which are either complex conjugates of each other or both real, such that for every positive integer $r$
  $$K(\chi^{(r)};a,b)=\sum_{\gm\in\bbF_{q^r}^*}\chi^{(r)}(a\gm+b\gm^{-1})=-w_1^r-w_2^r,$$
\end{Theorem}
\begin{proof}
  Let $G$ be the set of all rational functions $r(X)\in\bbF_q(X)$ such that for any $r\in G$, $r(X)=\nicefrac{h_1(X)}{h_2(X)}$ where $h_1,h_2\in\bbF_q[x]$ and both $h_1,h_2$ are monic. Let $\ovG$ the subgroup of $G$ such that for any $r\in \ovG$, $r(X)=\nicefrac{h_1(X)}{h_2(X)}$ where $h_1,h_2$ have non-zero constant term. We will define a multiplicative function $\lm:G\to \bbC$ such that $|\lm(r)|\leq 1$ for all $r\in G$ and $\lm(1)=1$. Let $r\in \ovG$ where $r(X)=\frac{h_1(X)}{h_2(X)}$ and in common splitting field $E$ of $h_1,h_2$ they factor as $$h_1(X)=\prod_{i=1}^{d_1}(X-\alpha_i)\qquad\text{and}\qquad h_2(X)=\prod_{j=1}^{d_2}(X-\beta_j)$$Then we will define $\lm$ using $\chi$ in the following way: \begin{equation}
    \lm(r)=\chi\lt[a\lt(\sum_{i=1}^{d_1}\alpha_i-\sum_{j=1}^{d_2}\beta_j  \rt) +b\lt( \sum_{i=1}^{d_1}\frac1{\alpha_i}-\sum_{j=1}^{d_2}\frac1{\beta_j} \rt)\rt]
  \end{equation}Then if $h\in\Phi_n$ such that $h(X)=X^n-c_1\cdot X^{n-1}+\cdots +(-1)^n\cdot c_n$ then $\lm(h)=\chi\lt(a\cdot c_1+b\cdot c_{k-1}\cdot c_k^{-1}\rt)$. Now it is enough to show that the corresponding coefficients are additive if two polynomials are multiplied since $\chi$ is an additive character. Let $h_1\in\Phi_{n_1}$ and $g_2\in \Phi_{n_2}$ and $$g_1(X)=X^{n_1}+\sum_{i=1}^{n_1}(-1)^i\cdot c_{1,i}\cdot X^{n_1-i},\quad g_2(X)=X^{n_2}+\sum_{j=2}^{n_2}(-2)^j\cdot c_{2,j}\cdot X^{n_2-j}\text{ and }g_1\cdot g_2=X^{n_1+n_2}+\sum_{t=1}^{n_1+n_2}(-1)^t\cdot c_j\cdot X^{n_1+n_2-t}$$Then by comparing coeffieicnets we have $$c_1=c_{1,1}+c_{2,1},\quad c_{n_1+n_2}=c_{1,n_1}\cdot c_{2,n_2},\quad\text{and}\quad c_{n_1+n_2-1}=c_{1,n_1}\cdot c_{2,n_2-1}+c_{1,n_1-1}\cdot c_{2,n_2}$$Hence $\lm$ is indeed a multiplicative function. Now we extend $\lm$ to $G$ by setting $\lm(r)=0$ for all $h\notin \ovG$. 

  Now for any $k\in\bbZ_0$,if  $(c_1,c_{k-1},c_k)$ is fixed where $c_k\neq 0$, there are $q^{k-3}$ possible choices for $c_2,\dots, c_{k-2}$ such that all polynomials with coefficients $c_i$'s give same value for $\chi(a\cdot c_1+b\cdot  c_{k-1}\cdot c_k^{-1}$. So we will show that for $k\geq 3$, $\sum_{h\in\Phi_k}\lm(h)=0$. So 
  \begin{align*}
    \sum_{g\in\Phi_k}\lm(g) & = \sum_{c_1,\dots, c_{k-1}}\sum_{c_k\in\bbF_q^*}\chi(a\cdot c_1+b\cdot c_{k-1}\cdot c_k^{-1})\\ 
    & = q^{k-3}\sum_{c_1,c_{k-1}\in \bbF_q}\sum_{c_k\in\bbF_q^*}\chi(a\cdot c_1+b\cdot c_{k-1}\cdot c_k^{-1})\\ 
    & = q^{k-3}\lt(\sum_{c_1\in\bbF_q}\chi(a\cdot c_1)\rt)\lt( \sum_{c_{k-1}\in \bbF_q}\sum_{c_k\in\bbF_q^*}\chi(b\cdot c_{k-1}\cdot c_k^{-1})\rt)= 0 \by{\AddCharPropsthm[(i)]}
  \end{align*} Hence we have $\ovL(z,\lm)=1+z\sum_{h\in\Phi_1}\lm(h)+z^2\sum_{h\in\Phi_2}\lm(h)$. Now $\Phi_1=\{X-\gm\mid \gm\in\bbF_q^*\}$. Therefore we have $$\sum_{h\in\Phi_1}\lm(h)=\sum_{\gm\in\bbF_q^*}\chi(a\cdot \gm+b\cdot \gm^{-1})=K(\chi;a,b)$$And furthermore for $\Phi_2$ we get \begin{align*}
    \sum_{h\in\Phi_2}\lm(h) & = \sum_{c_1\in\bbF_q}\sum_{c_2\in\bbF_q^*}\lm(X^2-c_1\cdot X+c_2)\\ 
    & = \sum_{c_1\in\bbF_q}\sum_{c_2\in\bbF_q^*}\chi(a\cdot c_1+b\cdot c_1\cdot c_2^{-1})\\ 
    & = \sum_{c_1\in\bbF_q}\sum_{c_2\in\bbF_q^*}\chi(c_1(a+b\cdot c_2^{-1}))
  \end{align*}
  So if $c_2=-a^{-1}b$ then $a+b\cdot c_2^{-1}=0$ and hence $\sum_{c_1\in\bbF_q}\chi(0)=q$ and if $c_2\neq -a^{-1}b$ then $\sum_{c_1\in\bbF_q}\chi(c_1(a+b\cdot c_2^{-1}))=0$ by \AddCharPropsthm[(i)]. So therefore $\sum_{h\in\Phi_2}\lm(h)=q$ and thus we get $\ovL(z,\lm)=1+k(\chi;a,b)\cdot z+q\cdot z^2$. Let $w_1,w_2\in\bbC$ which are either complex conjugates or both real such that $$\ovL(z,\lm)=(1-w_1\cdot z)(1-w_2\cdot z)$$

  Now we will calculate $L_r$  where $r=[E\colon \bbF_q]$. From the definition of $L_r$ in previous section we have $$L_r =  \sum_{\substack{h\in\irr(\Phi),\\ \deg(h)\mid r}}\deg(h)\cdot \chi(h)^{\nicefrac{r}{\deg(h)}}= \sum_{\substack{h\in\irr(\Phi),\\ \deg(h)\mid r}} \deg(h)\cdot \lm\lt( h^{\nicefrac{r}{\deg(h)}}\rt)$$Now for any irreducible $h$ let $\gm$ be a root of $h$ in the splitting field of $f$. Let $$h^{\nicefrac{r}{\deg(f)}}= x^r-c_1\cdot x^{r-1}+\cdots+(-1)^k c_r=\prod_{i=o}^{r-1}\lt(X-\gm^{q^i}\rt)$$Hence 
  $c_1=\tr_{E/\bbF_q}(\gm)$ and $c_r=\norm_{E/\bbF_q}(\gm)$. Hence we have $$\lm\lt(h^{\nicefrac{r}{\deg(h)}}\rt)=\psi(c_r)\cdot \chi(c_1)=\psi\lt(\norm_{E/\bbF_q}(\gm)\rt)\cdot \chi\lt(\tr_{E/\bbF_q}(\gm)\rt)=\psi^{(r)}(\gm)\cdot \chi^{(r)}(\gm)$$Now we can express $\deg(h)\cdot \lm\lt(h^{\nicefrac{r}{\deg(h)}}\rt)=\sum_{\gm\in E,h(\gm)=0}\psi^{(r)}(\gm)\cdot \chi^{(r)}(\gm)$. 
  Therefore $$L_r= \sum_{\substack{h\in\irr(\Phi),\\ \deg(h)\mid k}} \deg(h)\cdot \lm\lt( h^{\nicefrac{k}{\deg(h)}}\rt)=\sum_{\substack{h\in\irr(\Phi),\\ \deg(h)\mid k}}\sum_{\gm\in E,h(\gm)=0}\psi^{(r)}(\gm)\cdot \chi^{(r)}(\gm) $$Now all irreducible polynomials of $\bbF_q$ whose degree divides $r$ with non-zero constant their roots in $E$ are disjoint sets and partition $E^*$. Therefore we have $$L_r=\sum_{\gm\in E^*}\psi^{(r)}(\gm)\cdot \chi^{(r)}(\gm)=G(\psi^{(r)},\chi^{(r)})$$Since  by \LrootFactthm we have $L_r=-w_1^r-w_2^r$. So we have the theorem. 
\end{proof}

Assuming $w_1,w_2$ as above, every lifted Kloosterman sum is a symmetric function of $w_1,w_2$, which by Waring's identities can be rewritten purely in terms of $K=K(\chi;a,b)$ and $q$. This gives the following closed-form reduction of $K(\chi^{(s)};a,b)$ to the ground-field sum.
\begin{lemma}{}{thm:kloosterman-reduction}
  Let $\chi\in\Xq$ be a non-trivial additive character of $\bbF_q$ and let $a,b\in\bbF_q$ with $ab\neq 0$. Then for every positive integer $s$,
  $$K\lt(\chi^{(s)};a,b\rt)=\sum_{j=0}^{\lfloor s/2\rfloor}(-1)^{s-j-1}\cdot \frac{s}{s-j}\cdot \binom{s-j}{j}\cdot q^j\cdot K(\chi;a,b)^{s-2j}.$$
\end{lemma}
\begin{proof}
  Consider the polynomial $X_1^r+X_2^r$. Now by \Waringthm we can express this polynomial in terms of elementary symmetric polynomials $\esym_1,\dots, \esym_r$ which gives \begin{align*}
    X_1^r+X_2^r& = \sum_{\substack{d_1,d_2\in\bbZ_0\\ d_1+2d_2=r}}(-1)^{d_2}\;\frac{r(d_1+d_2-1)!}{d_1!d_2!}\esym_1(X_1,X_2)^{d_1}\cdot \esym_2(X_1,X_2)^{d_2}\\ 
    & = \sum_{j=0}^{\lfloor\nicefrac{r}{2}\rfloor}(-1)^j \;\frac{(r-j-1)!r}{j!(s-2j)!}\esym_1(X_1,X_2)^{s-2j}\cdot \esym_2(X_1,X_2)^j\\ 
    & = \sum_{j=0}^{\lfloor\nicefrac{r}{2}\rfloor}(-1)^j \binom{r-j}{j}\;\frac{r}{s-j}\esym_1(X_1,X_2)^{s-2j}\cdot \esym_2(X_1,X_2)^j
  \end{align*}
  Now we take $X_1=w_1$ and $X_2=w_2$ obtained from \thmref{thm:kloosterman-lift} we get $w_1^k+w_2^k=-K(\chi^{(k)};a,b)$ for any $k\in\bbN$ and $w_1\cdot w_2=q$. So by this replacing in the equation we get the result. 
\end{proof}
This lemma gives a relation between powers of Kloosterman sum and lifted Kloosterman sums. We can obtain a recursive relation between different lifts of Kloosterman sums very easily using this identiy$$w_1^r+w_2^r=(w_1^{r-1}+w_2^{r-1})(w_1+w_2)-(w_1^{r-2}+w_2^{r-2})w_1w_2$$ which gives us $$K(\chi^{(r)};a,b)=-K(\chi^{(r-1)};a,b)\cdot K(\chi;a,b)-q\cdot K(\chi^{(r-2)};a,b)$$ and so we have the corollary. 
\begin{corolary}{}{cor:kloosterman-recursion}
  Let $\chi\in\Xq$ be a non-trivial additive character of $\bbF_q$ and let $a,b\in\bbF_q$ with $ab\neq 0$. Then for every positive integer $s$, the lifted Kloosterman sums satisfy the recurrence $$K\lt(\chi^{(s)};a,b\rt)=-K\lt(\chi^{(s-1)};a,b\rt)\cdot K(\chi;a,b)-K\lt(\chi^{(s-2)};a,b\rt)\cdot q\qquad \text{for all }s\geq 2,$$ with the convention $K(\chi^{(0)};a,b)\coloneqq -2$ (and noting $K\lt(\chi^{(1)};a,b\rt)=K(\chi;a,b)$).
\end{corolary}
